\documentclass[UTF8,fontset=fandol,11pt]{ctexart}
\usepackage[a4paper,margin=24mm]{geometry}
\usepackage{amsmath,amssymb,amsthm,graphicx,tabularx,booktabs,hyperref}
\newtheorem{definition}{定义}[section]
\newtheorem{theorem}[definition]{定理}
\newtheorem{example}[definition]{例题}
\newtheorem{exercise}[definition]{练习}
\hypersetup{unicode=true,colorlinks=true,linkcolor=blue}
\title{数列极限 · 本科教学讲义}
\date{}
\begin{document}
\maketitle
\par\noindent 离线合成演示：不代表真实模型对本专题的教学验收。\par
\section{数列极限的定义}\label{sec:S01}
% UNIT S01.U01
\begin{definition}
称实数列 $a_n$ 收敛于 $a$，如果对每个 $\varepsilon>0$，存在正整数 $N$，使得所有 $n>N$ 都满足 $|a_n-a|<\varepsilon$。
\end{definition}
% UNIT S01.U02
\begin{theorem}
收敛实数列的极限唯一。
\end{theorem}
% UNIT S01.U03
\begin{proof}
若极限 $a\ne b$，取 $\varepsilon=|a-b|/3>0$。由两次收敛定义，充分大的 $n$ 同时满足 $|a_n-a|<\varepsilon$ 和 $|a_n-b|<\varepsilon$。三角不等式给出 $|a-b|<2|a-b|/3$，矛盾。
\end{proof}
% UNIT S01.U04
\begin{example}
证明 $a_n=1/n$ 收敛于 $0$。
\end{example}
% UNIT S01.U05
\par\noindent\textbf{解：}任取 $\varepsilon>0$，取正整数 $N>1/\varepsilon$。当 $n>N$ 时，$|1/n-0|=1/n<1/N<\varepsilon$。故所求极限为 $0$。\par
% UNIT S01.E01
\begin{exercise}
按定义证明 $3/n\to0$。
\end{exercise}
% UNIT S01.A01
\par\noindent\textbf{解：}取正整数 $N>3/\varepsilon$，则 $n>N$ 时 $3/n<\varepsilon$。\par
\section{极限的运算}\label{sec:S02}
% UNIT S02.U01
\begin{definition}
若 $a_n\to a$，则 $r_n=a_n-a$ 称为余项，并且 $r_n\to0$。
\end{definition}
% UNIT S02.U02
\begin{theorem}
若实数列 $a_n\to a$ 且 $b_n\to b$，则 $a_n+b_n\to a+b$。
\end{theorem}
% UNIT S02.U03
\begin{proof}
任取 $\varepsilon>0$，分别对两个收敛数列使用 $\varepsilon/2$。充分大的 $n$ 满足 $|(a_n+b_n)-(a+b)|\leq |a_n-a|+|b_n-b|<\varepsilon$。
\end{proof}
% UNIT S02.U04
\begin{example}
求 $\lim_{n\to\infty}(2+1/n)$。
\end{example}
% UNIT S02.U05
\par\noindent\textbf{解：}常数列 $2$ 的极限为 $2$，且 $1/n\to0$。由加法法则，所求极限为 $2$。\par
% UNIT S02.E01
\begin{exercise}
按定义证明 $5/n\to0$。
\end{exercise}
% UNIT S02.A01
\par\noindent\textbf{解：}取正整数 $N>5/\varepsilon$，则 $n>N$ 时 $5/n<\varepsilon$。\par
\end{document}
